\documentclass{article}%
\usepackage{float}%
\usepackage{hyperref}%
\usepackage{amssymb}%
\usepackage{amsthm}%
\usepackage{amsmath}%
\usepackage{stmaryrd}%
\usepackage{thm-restate}%
\usepackage{microtype}%
\usepackage{mathpartir}%
\usepackage{nameref}%
\usepackage{tikz-cd}%
\usepackage{todonotes}%
\usepackage{makecell}%
\usepackage{xspace}%
\usepackage{quiver}%
\usepackage{xcolor}%
\usepackage{fvextra}%
\usepackage{mdframed}%
\usepackage{forest}

\usepackage{ebproof}

\IfFileExists{tikzlibrarynfold.code.tex}{%
  \usetikzlibrary{nfold}%
  \tikzset{triple arrow/.style={scaling nfold=3}}%
}{%
  % tikz-nfold was added after TeX Live 2022.  Approximate its three-line
  % arrows so the document remains buildable on older installations.
  \tikzset{triple arrow/.style={%
    double distance between line centers=.9ex,
    postaction={draw, line width=.4pt, -}%
  }}%
}

\usepackage{macros}%

\definecolor{githubcodebackground}{HTML}{F6F8FA}%
\definecolor{githubcodetext}{HTML}{1F2328}%

\DefineVerbatimEnvironment{githubcode}{Verbatim}{%
  fontsize=\footnotesize,
  formatcom=\color{githubcodetext},
  breaklines=true,
  breakanywhere=true,
  breaksymbolleft={},
  vspace=0pt
}%
\surroundwithmdframed[%
  backgroundcolor=githubcodebackground,
  linecolor=githubcodebackground,
  linewidth=0pt,
  leftmargin=0pt,
  rightmargin=0pt,
  innerleftmargin=8pt,
  innerrightmargin=8pt,
  innertopmargin=7pt,
  innerbottommargin=7pt,
  splittopskip=7pt,
  splitbottomskip=7pt,
  skipabove=\medskipamount,
  skipbelow=\medskipamount
]{githubcode}%

\newtheorem{theorem}{Theorem}
\newtheorem{proposition}[theorem]{Proposition}
\newtheorem{lemma}[theorem]{Lemma}

\theoremstyle{remark}
\newtheorem{remark}[theorem]{Remark}

\theoremstyle{definition}
\newtheorem{example}[theorem]{Example}

\theoremstyle{definition}
\newtheorem{definition}[theorem]{Definition}

\theoremstyle{plain}
\newtheorem{corollary}[theorem]{Corollary}

%%%%%%%%%%%%%%%%%%%%%%%%%%%%%%%%%%%%%%%%%%%%%%%%%%%%%%%%%%%%%%%%%
%% Induction rules

\usepackage{mathpartir}
\renewcommand{\TirNameStyle}{\textsc}
% style of the induction rule name for mathpartir

%%%%%%%%%%%%%%%%%%%%%%%%%%%%%%%%%%%%%%%%%%%%%%%%%%%%%%%%%%%%%%%%%%%%%%%%%%%%%%%%%%%%%%
%Reference to inference rules, inspired by
% https://tex.stackexchange.com/questions/340788/cross-referencing-inference-rules

\makeatletter

\NewDocumentCommand \mpr@label { m }{%
  \def\@currentlabelname{#1}%
   \phantomsection % Creating an anchor to get a correct hyper reference link
}

\DeclareDocumentCommand \inferdef {m m m}{%
   \inferrule*[vcenter,right=#1]{#2}{#3}
   \mpr@label{\textsc{#1}}
}%

\makeatother

\newcommand{\ruleref}[1]{Rule \nameref{#1}}


\title{Program Length and Optimal Explicit Encodings of CaTT Entities}
\author{Karol Ołtarzewski}

\begin{document}

\maketitle

\begin{abstract}
We characterize straight-line programs whose instructions are operations derived
from the \catt syntactic construction rules. These programs construct syntax
trees of \catt entities. We introduce explicit encodings which measure the size
of a \catt entity both by the number of program instructions and by the encoded
bit-length of a program constructing it under a worst-case optimal explicit
encoding.
\end{abstract}

\section{Introduction}
Finster and Mimram have defined a dependent type theory called \catt, which
describes the structure of \(\omega\)-categories. The type theory \catt has four
kinds of syntactic entitites: contexts, substitutions, types, and terms. In
order to construct these entities one need up to eight syntactic instructions.
In this article we characterize straight-line programs whose instructions are
the eight elementary operations. Such a program constructs an abstract syntax
tree of a \catt entity, one vertex at a time. We introduce explicit encodings of
such programs. An explicit encoding assigns codewords to the operation tags and
to the parameters, using a pattern code with a common suffix terminator. This
yields two measures of size: the number of program instructions, which equals
the number of vertices of the constructed tree, and the encoded bit-length of
the program, which counts the total number of symbols in its explicit
presentation. The first measure is invariant to the encoding, while the second
depends on the assignment of codewords. For a finite family of \catt syntax
trees, we ask for an explicit encoding that minimizes the largest encoded
program length over the family, and we call such an encoding worst-case optimal.
We give a finite search that constructs a worst-case optimal explicit encoding
over a fixed alphabet, and we study basic properties of these encodings.

\paragraph{Related works.}
Yanofsky studied Kolmogorov complexity of categories and for that purpose
developed \emph{Sammy} language~\cite{Yanofsky}.Benjamin, Markakis, Offord,
Sarti, and Vicary measured the sizes of higher Eckmann--Hilton terms by counting
characters of a particular syntactic representation~\cite{Beyond_eh}. Benjamin
did similar measurements for higher identity proofs in homotopy type
theory~\cite{benjamin_higher_identity}.

\section{\catt syntax and elementary operations}
Let us briefly recall type theory \catt. Finster--Rice--Vicary provided a
compact set of rules that yields the abstract syntax construction of \catt
entities, the formation rules are:~\cite{ASfSAaUC}

\begin{figure}[H]
  \centering
  \small
  \begin{tabular}{@{}c@{\hspace{3em}}c@{}}
    \begin{prooftree}
      \infer0{\emptyset : \mathsf{Ctx}}
    \end{prooftree}
    &
    \begin{prooftree}
      \hypo{\Gamma : \mathsf{Ctx}}
      \hypo{A : \mathsf{Type}_\Gamma}
      \infer2{\Gamma, (x : A) : \mathsf{Ctx}}
    \end{prooftree}
    \\[2.2em]
    \begin{prooftree}
      \infer0{\langle \rangle : \emptyset \to \Gamma}
    \end{prooftree}
    &
    \begin{prooftree}
      \hypo{\sigma : \Delta \to \Gamma}
      \hypo{A : \mathsf{Type}_\Delta}
      \hypo{t : \mathsf{Term}_\Gamma}
      \infer3{
        \langle \sigma, x \mapsto t \rangle
        : (\Delta, x : A) \to \Gamma
      }
    \end{prooftree}
    \\[2.2em]
    \begin{prooftree}
      \infer0{\obj : \mathsf{Type}_\Gamma}
    \end{prooftree}
    &
    \begin{prooftree}
      \hypo{A : \mathsf{Type}_\Gamma}
      \hypo{s : \mathsf{Term}_\Gamma}
      \hypo{t : \mathsf{Term}_\Gamma}
      \infer3{\arr[A]{s}{t} : \mathsf{Type}_\Gamma}
    \end{prooftree}
    \\[2.2em]
    \begin{prooftree}
      \hypo{x \in \Gamma\vphantom{\mathsf{Type}}}
      \infer1{x : \mathsf{Term}_\Gamma}
    \end{prooftree}
    &
    \begin{prooftree}
      \hypo{\Delta : \mathsf{Ctx}}
      \hypo{A : \mathsf{Type}_\Delta}
      \hypo{\sigma : \Delta \to \Gamma}
      \infer3{\coh_{\Delta,A}[\sigma] : \mathsf{Term}_\Gamma}
    \end{prooftree}
  \end{tabular}
  \caption{Abstract syntax constructions in \catt.}
  \label{fig:syntax}
\end{figure}

\begin{figure}[H]
  \centering
  \small

  \begin{tabular}{@{}c@{\hspace{3em}}c@{}}
    \begin{prooftree}
      \infer0{\emptyset \vdash}
    \end{prooftree}
    &
    \begin{prooftree}
      \hypo{\Gamma \vdash}
      \hypo{\Gamma \vdash A}
      \infer2{\Gamma, (x : A) \vdash}
    \end{prooftree}
    \\[2.2em]
    \begin{prooftree}
      \infer0{\Gamma \vdash \obj}
    \end{prooftree}
    &
    \begin{prooftree}
      \hypo{\Gamma \vdash A}
      \hypo{\Gamma \vdash s : A}
      \hypo{\Gamma \vdash t : A}
      \infer3{\Gamma \vdash \arr[A]{s}{t}}
    \end{prooftree}
    \\[2.2em]
    \begin{prooftree}
      \infer0{\Delta \vdash \langle \rangle : \emptyset}
    \end{prooftree}
    &
    \begin{prooftree}
      \hypo{\Delta \vdash \sigma : \Gamma}
      \hypo{\Gamma \vdash A}
      \hypo{\Delta \vdash t : A\llbracket\sigma\rrbracket}
      \infer3{
        \Delta \vdash
        \langle \sigma, x \mapsto t \rangle
        : (\Gamma, x : A)
      }
    \end{prooftree}
    \\[2.2em]
    \multicolumn{2}{c}{
      \begin{prooftree}
        \hypo{\Gamma \vdash}
        \hypo{(x : A) \in \Gamma}
        \infer2{\Gamma \vdash x : A}
      \end{prooftree}
    }
  \end{tabular}

  \par\bigskip

  % Composition rule.
  \begin{prooftree}
    \hypo{
      \begin{gathered}
        \Delta \vdash_{\mathrm{p}}
        \qquad
        \Delta \vdash \arr[A]{s}{t}
        \qquad
        \Gamma \vdash \sigma : \Delta
        \\[0.4em]
        \operatorname{supp}(s)
        =
        \operatorname{supp}
          \bigl(\delta_{\Delta}^{-,\dim(\Delta)-1}\bigr)
        \qquad
        \operatorname{supp}(t)
        =
        \operatorname{supp}
          \bigl(\delta_{\Delta}^{+,\dim(\Delta)-1}\bigr)
      \end{gathered}
    }
    \infer1{
      \Gamma \vdash
      \coh_{\Delta,\arr[A]{s}{t}}[\sigma]
      :
      \arr[{A\llbracket\sigma\rrbracket}]
        {s\llbracket\sigma\rrbracket}
        {t\llbracket\sigma\rrbracket}
    }
  \end{prooftree}

  \par\bigskip

  % Coherence rule.
  \begin{prooftree}
    \hypo{
      \begin{gathered}
        \Delta \vdash_{\mathrm{p}}
        \qquad
        \Delta \vdash \arr[A]{s}{t}
        \qquad
        \Gamma \vdash \sigma : \Delta
        \\[0.4em]
        \operatorname{supp}(s)
        =
        \operatorname{supp}(t)
        =
        \Delta
      \end{gathered}
    }
    \infer1{
      \Gamma \vdash
      \coh_{\Delta,\arr[A]{s}{t}}[\sigma]
      :
      \arr[{A\llbracket\sigma\rrbracket}]
        {s\llbracket\sigma\rrbracket}
        {t\llbracket\sigma\rrbracket}
    }
  \end{prooftree}

  \par\medskip
  \caption{Typing rules in \catt.}
  \label{fig:typing}
\end{figure}

Rules above determine eight elementary operations. These operations, assembled
as straight-line program, that is, finite sequence of steps executed in order,
without loops, can construct every term, type, substitution or context of \catt.
Syntactic representations of these eight elementary operations are following:
% These programs are in normal form and are encodings of abstract syntax trees
% \catt construction. They exhibit quine-like property -- their source code is
% their output.

\begin{center}
\small
\setlength{\tabcolsep}{4pt}
\renewcommand{\arraystretch}{1.25}
\makebox[\textwidth][c]{%
\begin{tabular}{@{}lllc@{}}
\hline
Operation
  & Argument tree
  & Result
  & Corresponding rule from Figure~\ref{fig:syntax} \\
\hline

$\mathtt{NODE}(\mathsf{EmptyCtx})$
  & none
  & $\emptyset$
  & $\dfrac{}{\emptyset : \mathsf{Ctx}}$
  \\[1.4em]

$\mathtt{NODE}(\mathsf{ExtendCtx},x)$
  & $\Gamma,A$
  & $(\Gamma,x:A)$
  & $\dfrac{
        \Gamma : \mathsf{Ctx}
        \quad A : \mathsf{Type}_\Gamma
      }{
        \Gamma,(x:A) : \mathsf{Ctx}
      }$
  \\[1.4em]

$\mathtt{NODE}(\mathsf{EmptySub})$
  & none
  & $\langle\rangle$
  & $\dfrac{}{
        \langle\rangle : \emptyset \to \Gamma
      }$
  \\[1.4em]

$\mathtt{NODE}(\mathsf{ExtendSub},x)$
  & $\sigma,t$
  & $\langle\sigma,x\mapsto t\rangle$
  & $\dfrac{
        \sigma : \Delta \to \Gamma
        \quad A : \mathsf{Type}_\Delta
        \quad t : \mathsf{Term}_\Gamma
      }{
        \langle\sigma,x\mapsto t\rangle
        : (\Delta,x:A) \to \Gamma
      }$
  \\[1.4em]

$\mathtt{NODE}(\mathsf{Star})$
  & none
  & $\obj$
  & $\dfrac{}{\obj : \mathsf{Type}_\Gamma}$
  \\[1.4em]

$\mathtt{NODE}(\mathsf{Arrow})$
  & $A,s,t$
  & $\arr[A]{s}{t}$
  & $\dfrac{
        A : \mathsf{Type}_\Gamma
        \quad s : \mathsf{Term}_\Gamma
        \quad t : \mathsf{Term}_\Gamma
      }{
        \arr[A]{s}{t} : \mathsf{Type}_\Gamma
      }$
  \\[1.4em]

$\mathtt{NODE}(\mathsf{Var},x)$
  & none
  & $x$
  & $\dfrac{x \in \Gamma}{x : \mathsf{Term}_\Gamma}$
  \\[1.4em]

$\mathtt{NODE}(\mathsf{Coh})$
  & $\Delta,A,\sigma$
  & $\coh_{\Delta,A}[\sigma]$
  & $\dfrac{
        \Delta : \mathsf{Ctx}
        \quad A : \mathsf{Type}_\Delta
        \quad \sigma : \Delta \to \Gamma
      }{
        \coh_{\Delta,A}[\sigma] : \mathsf{Term}_\Gamma
      }$
  \\[.8em]

\hline
\end{tabular}%
}
\caption{\catt elementary operations.}
\label{tab:elementary-operations}
\end{center}

Each of these eight operations has a unique tag: \texttt{EmptyCtx},
\texttt{ExtendCtx}, \texttt{EmptySub}, \texttt{ExtendSub}, \texttt{Star},
\texttt{Arrow}, \texttt{Var}, or \texttt{Coh}. The operations
\texttt{ExtendCtx}, \texttt{ExtendSub}, and \texttt{Var} have an additional
parameter, specifying name of the term being introduced. Every operation can be
therefore specified by at most two fields, tag field and parameter field. The
operations whose tags are \texttt{ExtendCtx}, \texttt{ExtendSub},
\texttt{Arrow}, and \texttt{Coh} also require argument trees, as listed in the
table above.

\section{Abstract Syntax Trees and Construction Programs}
We will denote \(\mathcal{T} \ni T\) as family of all \catt abstract syntax
trees derived from syntactic construction build according to
rules~\ref{fig:syntax} and ~\ref{fig:typing}. \(\mathcal{T}_\text{fin} \subseteq
\mathcal{T}\) is family of finite trees. Vertices of \(T\) correspond to
operations from \ref{tab:elementary-operations}. \(T\) can
represent any \catt entity -- context, substitution, type, or term.

\begin{remark}
Different typing judgments can have the same abstract syntax tree for their
syntactic entity. To see that different typing judgments can have the same tree,
consider two typing judgments
\[
  (x:\obj)\vdash x:\obj,
  \qquad
  (x:\obj,y:\obj)\vdash x:\obj.
\]
Both are derivable using Figure~\ref{fig:typing}. Their syntactic
entities have the same abstract syntax tree, consisting of a single vertex
labelled \(\mathsf{Var}\) with parameter \(x\).
\end{remark}


% \begin{definition}
Let \(E\) be a finite encoding, over alphabet \(\Sigma_E\), of each of
elementary \catt operations. Encodings are independent of arbitrary notation
used to represent these operations. \(\mathcal{P}_E\) is a function which
assigns each \catt abstract syntax tree a family of sequences of \catt
elementary operations in a given encoding \(E\), which describe intended syntax
tree. Each of these sequences is a program whose operations are ordered in a
manner permitted by rules~\ref{fig:syntax} and ~\ref{fig:typing}. In particular
\(\mathcal{P}_E(T)\) is a family of sequences where each sequence describe
\(T\in\mathcal{T}\). An element \(P \in \mathcal{P}_E(T)\) is an
\emph{admissible program}. These programs construct \(T\), starting with no
vertices. Every admissible program contains exactly one elementary operation for
each vertex of \(T\).
% \end{definition}

\begin{definition}[Length]
\label{def:program-length}
Let \(P=(o_1,\ldots,o_n)\in\mathcal{P}_E(T)\).
The \emph{length} of \(P\) is
\[
  |P| := n.
\]

For a fixed \(T\in\mathcal{T}\), every admissible program contains exactly one
elementary operation for each vertex of \(T\). Consequently, all programs in
\(\mathcal{P}_E(T)\) have the same length.
\end{definition}

Every \(T\in\mathcal{T}\) admits an elementary construction program: recursively
construct its argument subtrees, then apply the constructor at its root.
Consequently, \(\mathcal{P}_E(T)\) is nonempty.

Since every admissible program contains exactly one elementary operation for
each vertex occurrence of \(T\), all elements of \(\mathcal{P}_E(T)\) have the
same length:
\[
  |P| = |V(T)|
  \qquad\text{for every } P\in\mathcal{P}_E(T).
\]

\section{Structural Properties of Term Trees}
\begin{lemma}[\catt term tag counts]
\label{lem:catt-term-tag-counts}
Let \(T\in\mathcal{T}_\text{fin}\) be tree of a term, and let \(N_X\) denote the number
of occurrences of tag \(X\) in \(T\).
\begin{enumerate}
  \item \(N_{\mathtt{Coh}} = N_{\mathtt{EmptyCtx}} = N_{\mathtt{EmptySub}},\)
  \item \(N_{\mathtt{ExtendCtx}} = N_{\mathtt{ExtendSub}},\)
  \item \(N_{\mathtt{Star}} = N_{\mathtt{Coh}}+N_{\mathtt{ExtendCtx}},\) 
  \item \(N_{\mathtt{Var}} =
  2N_{\mathtt{Arrow}}+N_{\mathtt{ExtendCtx}}-N_{\mathtt{Coh}}+1. \)
  \item \(N_{\mathtt{ExtendCtx}} =
  N_{\mathtt{ExtendSub}} \ge N_{\mathtt{Coh}}\)
  \item \(N_{\mathtt{Arrow}}\ge N_{\mathtt{Coh}}\)
  \item \(N_{\mathtt{Var}} > N_{\mathtt{Arrow}}\)
  \item \(
  N_{\mathtt{Var}} > N_{\mathtt{Star}} > N_{\mathtt{ExtendCtx}} =
  N_{\mathtt{ExtendSub}} \ge N_{\mathtt{Coh}} = N_{\mathtt{EmptyCtx}} =
  N_{\mathtt{EmptySub}}
  \)
\end{enumerate}
\end{lemma}

\begin{proof}
\begin{enumerate}
\item[1. 2.] Each coherence contains context and substitution. Each context ends
with \(\mathtt{EmptyCtx}\), and each substitution ends with
\(\mathtt{EmptySub}\). Context and corresponding substitution have the same
number of elements. Summing over all coherences gives the first two identities.

\item[3.] Each context extension has one underlying type, each arrow has one
underlying type, and each coherence has one underlying type. All these positions
are \(\mathtt{Star}\) or \(\mathtt{Arrow}\), so
\[
  N_{\mathtt{Star}}+N_{\mathtt{Arrow}}=N_{\mathtt{ExtendCtx}}+N_{\mathtt{Arrow}}+N_{\mathtt{Coh}}.
\]

\item[4.] There is the root, there are two endpoints per arrow, and there is
one term per substitution extension. Each such position is occupied by
\(\mathtt{Var}\) or \(\mathtt{Coh}\), so
\[
  N_{\mathtt{Var}}+N_{\mathtt{Coh}}=1+2N_{\mathtt{Arrow}}+N_{\mathtt{ExtendCtx}}.
\]
\item[5. 6.] Each coherence has a context and an arrow type. Hence
\(N_{\mathtt{ExtendCtx}}\ge N_{\mathtt{Coh}}\) and \(N_{\mathtt{Arrow}}\ge
N_{\mathtt{Coh}}\).
\item[7.] By (4), (5) and (6). 
\item[8.]
\[
  N_{\mathtt{Arrow}} \geq N_{\mathtt{Coh}},
\]
and therefore
\[
  N_{\mathtt{Arrow}}+\frac12>N_{\mathtt{Coh}}.
\]
Multiplying by two and then adding
\(N_{\mathtt{ExtendCtx}}\) gives
\[
  2N_{\mathtt{Arrow}}+1+N_{\mathtt{ExtendCtx}}
  >
  2N_{\mathtt{Coh}}+N_{\mathtt{ExtendCtx}}.
\]
Subtracting \(N_{\mathtt{Coh}}\) from both sides, we obtain
\[
  2N_{\mathtt{Arrow}}+1+N_{\mathtt{ExtendCtx}}
  -N_{\mathtt{Coh}}
  >
  N_{\mathtt{Coh}}+N_{\mathtt{ExtendCtx}}.
\]
By identities (3) and (4), this is precisely
\[
  N_{\mathtt{Var}}>N_{\mathtt{Star}}.
\]
Other inequalities follow from previous proofs.
\end{enumerate}
\end{proof}

\section{Construction Orders and Examples}
\begin{proposition}[Admissible construction orders]
\label{prop:construction-orders}
Let \(T\in\mathcal{T}\), and choose an enumeration of its vertex occurrences.
Associate each occurrence with its elementary operation, using the argument
subtrees designated by \(T\) in their prescribed roles.

These operations can be performed in the chosen order to construct \(T\) from
the bottom up if and only if every child precedes its parent.
\end{proposition}

\begin{proof}
In a bottom-up construction, each argument subtree must be available before the
operation constructing its parent is performed. Hence every child must precede
its parent.

Conversely, suppose that every child precedes its parent. Proceed along the
enumeration. By induction, when a vertex is reached, its argument subtrees have
already been constructed. They remain available because each has a unique
parent, which has not yet been constructed. Applying the corresponding
elementary operation to these subtrees in their prescribed roles constructs the
subtree rooted at that vertex. The root occurs last, so the resulting tree is
\(T\).
\end{proof}

Proposition~\ref{prop:construction-orders} allows different construction
conventions, including left-to-right postorder, right-to-left postorder, and
interleaving independent subtree constructions, provided that every child is
constructed before its parent and prescribed roles are preserved. For the
example programs that follow, we use postorder traversal with stack execution:
argument trees are constructed in the order listed in the table, followed by the
operation constructing their parent. Each operation consumes its argument trees
from the top of the stack and pushes the resulting tree.

\begin{lemma}
Let \(P\) be a postorder program which constructs \catt coherence. \(P\) can be
split into a sequence of three subprograms, and one coherence operation, which
respectively construct the context, the type and the substitution.
\end{lemma}
\begin{proof}
The root of the syntax tree is a \texttt{Coh} operation, whose three ordered
argument subtrees represent the context, the type, and the substitution,
respectively.

By the definition of postorder traversal, each argument subtree is traversed
completely before the next one, and the root is visited last.
\end{proof}

In the example programs below, semicolons separate operations. Blank lines and
new lines are added only for readability.

\begin{example}[Binary composition]
\label{ex:comp}
Consider the context
\[
  \Delta_{\mathrm{c}}
  =
  \bigl(
    x:\obj,\,
    y:\obj,\,
    f:\arr[\obj]{x}{y},\,
    z:\obj,\,
    g:\arr[\obj]{y}{z}
  \bigr).
\]
The binary composite in this context is the term
\[
  \operatorname{comp}(f,g)
  =
  \coh_{\Delta_{\mathrm{c}},\,\arr[\obj]{x}{z}}
      [\mathrm{id}_{\Delta_{\mathrm{c}}}].
\]
Its syntax tree is constructed by the following program. The four blocks
construct the context \(\Delta_{\mathrm{c}}\), the type \(\arr[\obj]{x}{z}\),
the identity substitution, and the final coherence node, respectively.

\begin{githubcode}
NODE(EmptyCtx);
NODE(Star); NODE(ExtendCtx, x);
NODE(Star); NODE(ExtendCtx, y);
NODE(Star); NODE(Var, x); NODE(Var, y); NODE(Arrow);
NODE(ExtendCtx, f);
NODE(Star); NODE(ExtendCtx, z);
NODE(Star); NODE(Var, y); NODE(Var, z); NODE(Arrow);
NODE(ExtendCtx, g);

NODE(Star); NODE(Var, x); NODE(Var, z); NODE(Arrow);

NODE(EmptySub);
NODE(Var, x); NODE(ExtendSub, x);
NODE(Var, y); NODE(ExtendSub, y);
NODE(Var, f); NODE(ExtendSub, f);
NODE(Var, z); NODE(ExtendSub, z);
NODE(Var, g); NODE(ExtendSub, g);

NODE(Coh)
\end{githubcode}

Figure~\ref{fig:comp-syntax-tree} shows the resulting abstract syntax tree. Its
root is the \texttt{Coh} constructor. The three subtrees rooted at vertices
\(17\), \(21\), and \(32\) represent the context \(\Delta_{\mathrm{c}}\), the
type \(\arr[\obj]{x}{z}\), and the identity substitution, respectively. Children
are ordered from left to right. The gray bracketed numbers indicate the
positions of the corresponding operations in the postorder program. Variable
names are parameters of vertices.

\begin{figure}[H]
  \centering
  \begin{forest}
  for tree={
    grow=270,
    parent anchor=south,
    child anchor=north,
    anchor=north,
    align=center,
    font=\footnotesize,
    draw=black!45,
    rounded corners=2pt,
    fill=black!3,
    inner sep=1.5pt,
    l sep=8pt,
    s sep=2pt,
    edge={draw=black!55, line width=.35pt}
  }
    [{\compastlabel{33}{Coh}{}}
      [{\compastlabel{17}{ExtendCtx}{(g)}}
        [{\compastlabel{12}{ExtendCtx}{(z)}}
          [{\compastlabel{10}{ExtendCtx}{(f)}}
            [{\compastlabel{5}{ExtendCtx}{(y)}}
              [{\compastlabel{3}{ExtendCtx}{(x)}}
                [{\compastlabel{1}{EmptyCtx}{}}]
                [{\compastlabel{2}{Star}{}}]
              ]
              [{\compastlabel{4}{Star}{}}]
            ]
            [{\compastlabel{9}{Arrow}{}}
              [{\compastlabel{6}{Star}{}}]
              [{\compastlabel{7}{Var}{(x)}}]
              [{\compastlabel{8}{Var}{(y)}}]
            ]
          ]
          [{\compastlabel{11}{Star}{}}]
        ]
        [{\compastlabel{16}{Arrow}{}}
          [{\compastlabel{13}{Star}{}}]
          [{\compastlabel{14}{Var}{(y)}}]
          [{\compastlabel{15}{Var}{(z)}}]
        ]
      ]
      [{\compastlabel{21}{Arrow}{}}
        [{\compastlabel{18}{Star}{}}]
        [{\compastlabel{19}{Var}{(x)}}]
        [{\compastlabel{20}{Var}{(z)}}]
      ]
      [{\compastlabel{32}{ExtendSub}{(g)}}
        [{\compastlabel{30}{ExtendSub}{(z)}}
          [{\compastlabel{28}{ExtendSub}{(f)}}
            [{\compastlabel{26}{ExtendSub}{(y)}}
              [{\compastlabel{24}{ExtendSub}{(x)}}
                [{\compastlabel{22}{EmptySub}{}}]
                [{\compastlabel{23}{Var}{(x)}}]
              ]
              [{\compastlabel{25}{Var}{(y)}}]
            ]
            [{\compastlabel{27}{Var}{(f)}}]
          ]
          [{\compastlabel{29}{Var}{(z)}}]
        ]
        [{\compastlabel{31}{Var}{(g)}}]
      ]
    ]
  \end{forest}
  \caption{The abstract syntax tree constructed by the program in
    Example~\ref{ex:comp}. Each of its \(33\) vertices corresponds to one
    \texttt{NODE} operation.}
  \label{fig:comp-syntax-tree}
\end{figure}

This program contains \(33\) elementary operations. Its output corresponds to
the right-hand side of the following definition in Benjamin's
syntax~\cite{Benjamin_Catt_Github}:

\begin{githubcode}
let comp_example (x(f)y(g)z) : x -> z = comp f g
\end{githubcode}
\end{example}

\begin{example}[Right whiskering]
\label{ex:whiskr}
Consider right whiskering in the context
\[
  \begin{aligned}
    \Delta_{\mathrm{w}} = \bigl(
      &x:\obj,\,
       y:\obj,\,
       f:\arr[\obj]{x}{y},\,
       g:\arr[\obj]{x}{y},\\
      &a:\arr[{\arr[\obj]{x}{y}}]{f}{g},\,
       z:\obj,\,
       h:\arr[\obj]{y}{z}
    \bigr).
  \end{aligned}
\]
Writing
\[
  A_{\mathrm{w}}
  =
  \arr[{\arr[\obj]{x}{z}}]
      {\operatorname{comp}(f,h)}
      {\operatorname{comp}(g,h)},
\]
the right-whiskering term is
\[
  \operatorname{whiskr}(a,h)
  =
  \coh_{\Delta_{\mathrm{w}},A_{\mathrm{w}}}
      [\mathrm{id}_{\Delta_{\mathrm{w}}}].
\]
The following program constructs its syntax tree.

\begin{githubcode}
NODE(EmptyCtx);
NODE(Star); NODE(ExtendCtx, x);
NODE(Star); NODE(ExtendCtx, y);
NODE(Star); NODE(Var, x); NODE(Var, y); NODE(Arrow);
NODE(ExtendCtx, f);
NODE(Star); NODE(Var, x); NODE(Var, y); NODE(Arrow);
NODE(ExtendCtx, g);
NODE(Star); NODE(Var, x); NODE(Var, y); NODE(Arrow);
NODE(Var, f); NODE(Var, g); NODE(Arrow);
NODE(ExtendCtx, a);
NODE(Star); NODE(ExtendCtx, z);
NODE(Star); NODE(Var, y); NODE(Var, z); NODE(Arrow);
NODE(ExtendCtx, h);

NODE(Star); NODE(Var, x); NODE(Var, z); NODE(Arrow);

NODE(EmptyCtx);
NODE(Star); NODE(ExtendCtx, x);
NODE(Star); NODE(ExtendCtx, y);
NODE(Star); NODE(Var, x); NODE(Var, y); NODE(Arrow);
NODE(ExtendCtx, f);
NODE(Star); NODE(ExtendCtx, z);
NODE(Star); NODE(Var, y); NODE(Var, z); NODE(Arrow);
NODE(ExtendCtx, g);
NODE(Star); NODE(Var, x); NODE(Var, z); NODE(Arrow);
NODE(EmptySub);
NODE(Var, x); NODE(ExtendSub, x);
NODE(Var, y); NODE(ExtendSub, y);
NODE(Var, f); NODE(ExtendSub, f);
NODE(Var, z); NODE(ExtendSub, z);
NODE(Var, h); NODE(ExtendSub, g);
NODE(Coh);

NODE(EmptyCtx);
NODE(Star); NODE(ExtendCtx, x);
NODE(Star); NODE(ExtendCtx, y);
NODE(Star); NODE(Var, x); NODE(Var, y); NODE(Arrow);
NODE(ExtendCtx, f);
NODE(Star); NODE(ExtendCtx, z);
NODE(Star); NODE(Var, y); NODE(Var, z); NODE(Arrow);
NODE(ExtendCtx, g);
NODE(Star); NODE(Var, x); NODE(Var, z); NODE(Arrow);
NODE(EmptySub);
NODE(Var, x); NODE(ExtendSub, x);
NODE(Var, y); NODE(ExtendSub, y);
NODE(Var, g); NODE(ExtendSub, f);
NODE(Var, z); NODE(ExtendSub, z);
NODE(Var, h); NODE(ExtendSub, g);
NODE(Coh);

NODE(Arrow);

NODE(EmptySub);
NODE(Var, x); NODE(ExtendSub, x);
NODE(Var, y); NODE(ExtendSub, y);
NODE(Var, f); NODE(ExtendSub, f);
NODE(Var, g); NODE(ExtendSub, g);
NODE(Var, a); NODE(ExtendSub, a);
NODE(Var, z); NODE(ExtendSub, z);
NODE(Var, h); NODE(ExtendSub, h);

NODE(Coh)
\end{githubcode}

This program contains \(117\) elementary operations. In Benjamin's syntax, the
coherence can be declared as follows:

\begin{githubcode}
coh whiskr (x(f(a)g)y(h)z) : comp f h -> comp g h
\end{githubcode}

The term constructed by the program is the right-hand side of the following
definition:

\begin{githubcode}
let whiskr_example (x(f(a)g)y(h)z)
    : comp f h -> comp g h
    = whiskr a h
\end{githubcode}
\end{example}

\section{Explicit Program Encodings}
Now, we approach complexity of \(P\) as a total number of symbols in its
explicit encoded presentation.

\begin{definition}[Total length]
\label{def:program-total-length}
Let \(P=(o_1,\ldots,o_n)\in\mathcal P_E(T)\), where each \(o_i\) is a complete
encoded operation over an alphabet. The \emph{total length} of \(P\) is
\[
  \|P\| := \sum_{i=1}^{n}|o_i|,
\]
where \(|o_i|\) denotes the number of symbols in \(o_i\). Equivalently,
\(\|P\|\) is the length of the string \(o_1\cdots o_n\) obtained by
concatenating operations.

For a fixed encoding \(E\) and tree \(T\), every program in \(\mathcal P_E(T)\)
contains the same encoded instructions, although their order may differ.
Consequently, all these programs have the same total length. We denote this
common value by
\[
  \|\mathcal P_E(T)\| := \|P\|,
  \qquad P\in\mathcal P_E(T).
\]
\end{definition}

\begin{proposition}
Let \(E\) use a fixed alphabet \(\Sigma_E\) containing \(q\geq 2\) symbols.
There exists a binary encoding \(E_{\mathrm{bin}}\) such that, for every
\(T\in\mathcal{T}_\text{fin}\) and every \(P\in\mathcal P_E(T)\), the corresponding
program \(P_{\mathrm{bin}}\in\mathcal P_{E_{\mathrm{bin}}}(T)\) satisfies
\[
  |P_{\mathrm{bin}}|=|P|,
  \qquad
  \|P_{\mathrm{bin}}\|
  =\lceil\log_2 q\rceil\,\|P\|.
\]
\end{proposition}

\begin{proof}
Apply the fixed-length alphabet encoding described in
\cite[Section~2.5]{greenlaw_hoover_1998}. Put \(b=\lceil\log_2 q\rceil\). Since
\(q\leq 2^b\), each symbol of \(\Sigma_E\) can be assigned a distinct binary
word of length \(b\). Define \(E_{\mathrm{bin}}\) by replacing every symbol in
each complete operation record with its assigned binary word.

The original symbols can be recovered by reading blocks of \(b\) bits. Applying
the decoder for \(E\) then recovers the original operation sequence. Therefore,
this replacement preserves the sequence of elementary operations, so
\(P_{\mathrm{bin}}\) describes the same tree \(T\) and satisfies
\(|P_{\mathrm{bin}}|=|P|\). Every encoded symbol is replaced by exactly \(b\)
bits, giving \(\|P_{\mathrm{bin}}\|=b\,\|P\|\).
\end{proof}

\begin{definition}[Worst-case optimal encoding]
\label{def:worst-case-optimal-encoding}
Let \(\mathcal F\subseteq\mathcal{T}_\text{fin}\), and let \(\mathfrak P\) be a
nonempty family of functions \(\mathcal P_E\).

An encoding \(E_*\), with
\(\mathcal P_{E_*}\in\mathfrak P\), is
\emph{worst-case optimal for \(\mathcal F\)} if
\[
  \max_{T\in\mathcal F}
    \|\mathcal P_{E_*}(T)\|
  =
  \min_{\mathcal P_E\in\mathfrak P}
    \max_{T\in\mathcal F}
      \|\mathcal P_E(T)\|.
\]
Thus, a worst-case optimal encoding minimizes the largest total program length
among the trees in \(\mathcal F\).
\end{definition}

\begin{definition}[Explicit program]
An encoded program \(P=(o_1,\ldots,o_n)\) is \emph{explicit} if each operation
\(o_i\) has all its fields separately encoded by codewords with a common fixed
nonempty suffix terminator, which cannot occur elsewhere within a codeword.
Equivalently, the field codewords form a pattern
code~\cite{anisimov_zavadskyi_2017}.
\end{definition}

\begin{remark}
No additional operation terminators or field separators are needed, since the
end of each field codeword is unambiguous. The operation tag field determines
whether a parameter field follows. Consequently, the concatenation \(o_1\cdots
o_n\) can be parsed as string unambiguously.
\end{remark}

\section{Worst-Case Optimal Explicit Encodings}
Now, we characterize a general construction for worst-case optimal explicit
encoding for a given nonempty finite family \(\mathcal
F\subseteq\mathcal{T}_\text{fin}\) over a fixed finite alphabet \(\Sigma\) with at
least two symbols.

Let \(\mathsf{Tag}\) be subset of eight tag fields, and let \(\mathsf{Param}\)
be a set of the parameter fields. We use code, that is, an injective function:
\[
  c\colon \mathsf{Tag}\sqcup\mathsf{Param}\longrightarrow\Sigma^+.
\]
whose codewords ending is a common suffix \(\delta\in\Sigma^+\) that occurs
nowhere else within any codeword. The disjoint union keeps tags distinct from
parameter indices. Each operation is encoded as
\[
  \begin{aligned}
    E(a)   &= c(a)     && \text{if \(a\) takes no parameter},\\
    E(a,p) &= c(a)c(p) && \text{if \(a\) takes a parameter \(p\in\mathsf{Param}\)}.
  \end{aligned}
\]
\(c(a)c(p)\) denotes the concatenation of tag and parameter codewords.

\begin{proposition}[Construction of an optimal explicit encoding]
\label{prop:constructive-minimax-encoding}
For every nonempty finite family \(\mathcal F\subseteq\mathcal{T}_\text{fin}\), there
is a finite search which finds an optimal encoding over alphabet \(\Sigma\) with
at least two symbols.
\end{proposition}
\begin{proof}
Let \(D=\mathsf{Tag}\sqcup\mathsf{Param}\) be the set of fields occurring in
\(\mathcal F\), and put \(d=|D|\). For each \(\delta\in\Sigma^+\), let
\[
  C_\delta := \{w\in\Sigma^+ :
    \text{\(\delta\) occurs exactly once in \(w\), as its suffix}\}.
\]
Perform the following search.
\begin{enumerate}
\item Enumerate the terminators \(\delta\in\Sigma^+\) of length at most \(d\).
\item For each \(\delta\), take the shortest \(d\) words of \(C_\delta\) in the
      fixed enumeration by nondecreasing length.
\item Try every bijection from \(D\) to these words and return a terminator and
      assignment minimizing \(\max_{T\in\mathcal F}\|\mathcal P_E(T)\|\).
\end{enumerate}

The search is finite. There are finitely many terminators in Step~1. For Step~2,
choose \(b\in\Sigma\) different from the first symbol of \(\delta\). The words
\(b^j\delta\), for \(0\leq j<d\), belong to \(C_\delta\), so it suffices to
inspect words of length at most \(|\delta|+d-1\). Step~3 tests \(d!\)
assignments per terminator on finitely many finite trees.
\end{proof}

\begin{remark}[Size of the search]
There are
\[
  \sum_{m=1}^{d}q^m=\frac{q^{d+1}-q}{q-1}
\]
suffixes in step~1 and at most
\[
  \frac{q^{d+1}-q}{q-1}\,d!
\]
assignments to evaluate, where \(q\) is number of symbols in prescribed
alphabet. For the binary alphabet this becomes \((2^{d+1}-2)d!\). Generating the
short codewords for each terminator require additional work.
\end{remark}


\begin{remark}[Predefined suffix]
If a particular suffix \(\delta\) is prescribed in advance, simply omit step~1.
The bound on used codeword lengths is then \(|\delta|+d-1\).
\end{remark}

\begin{proposition}[Enlarging the family]
\label{prop:minimax-family-monotonicity}
For nonempty finite families \(\mathcal F\subseteq\mathcal
G\subseteq\mathcal{T}_\text{fin}\), the optimal worst-case total length for \(\mathcal F\)
is at most that for \(\mathcal G\).
\end{proposition}
\begin{proof}
For every fixed encoding, the maximum over \(\mathcal F\) is at most the maximum
over \(\mathcal G\).
\end{proof}

\begin{proposition}[No single optimum for every finite family]
\label{prop:no-uniform-pattern-optimum}
No encoding \(E\) is worst-case optimal for every finite family of \catt
abstract syntax trees.
\end{proposition}
\begin{proof}
The one-vertex trees \(\texttt{EmptyCtx}\) and \(\texttt{Star}\) each contain
one tag field and no parameter. For either singleton family, the optimum is one
symbol assigned to the field. An encoding optimal for both singleton families
would give distinct one-symbol codewords to the two tags. However, two distinct
one-symbol words cannot have a common nonempty suffix. This contradicts the
pattern-code requirement.
\end{proof}

\begin{proposition}
Let \(\mathcal F_1\subseteq\mathcal{T}_\text{fin}\) and \(\mathcal
F_2\subseteq\mathcal{T}_\text{fin}\) be nonempty finite families such that the same
unique fields occur in both families. Within class of encodings with a fixed
terminator suffix, these families admit worst-case optimal encodings assigning
the same set of codewords to their occurring fields, differing only in the
assignment of these codewords.
\end{proposition}
\begin{proof}
Let \(N\) be the common number of unique fields. Choose \(N\) shortest codewords
permitted by the fixed terminator and list them in nondecreasing order of
length.

For either family, consider any admissible encoding and order its \(N\) used
codewords by nondecreasing length. Replace the \(j\)-th codeword by the \(j\)-th
chosen codeword. No field becomes longer, so no program in the family becomes
longer. Thus its worst-case total length cannot increase.

It therefore suffices to consider assignments of the chosen codewords to the
occurring fields. There are only finitely many such assignments, so one attains
the minimum. Applying this argument to both families proves the claim.
\end{proof}

\begin{lemma}[Frequency ordering]
\label{lem:frequency-ordering}
Every nonempty finite family \(\mathcal
F\subseteq\mathcal{T}_\text{fin}\) admits a worst-case optimal encoding with the following
property: if \(x\) occurs more often than \(y\) in some tree and less often in
none, then \(x\)'s codeword is no longer than \(y\)'s.
\end{lemma}
\begin{proof}
From the worst-case optimal encodings, choose one that minimizes the sum of the
encoded lengths of the trees in \(\mathcal F\). If \(x\) had the longer
codeword, interchanging the two would shorten at least one tree without
lengthening any other, contradicting this choice.
\end{proof}

\begin{proposition}
Let \(\mathcal F\subseteq\mathcal{T}_\text{fin}\) be a nonempty finite family of coherences,
and let \(c\) be the code of a worst-case optimal encoding, then
\[
  |c(\mathtt{Var})| \leq |c(\mathtt{Star})| \leq |c(\mathtt{ExtendCtx})|
\]
and
\[
  |c(\mathtt{Var})| \leq |c(\mathtt{Arrow})|
\]
Moreover, for \(\mathcal F\) there exists a worst-case optimal encoding whose
code \(\tilde{c}\) has the following properties:
\[
  |\tilde{c}(\mathtt{ExtendCtx})|  \leq |\tilde{c}(\mathtt{ExtendSub})| \leq 
  |\tilde{c}(\mathtt{Coh})| \leq |\tilde{c}(\mathtt{EmptyCtx})| \leq 
  |\tilde{c}(\mathtt{EmptySub})|
\]
and
\[
  |\tilde{c}(\mathtt{Arrow})| \leq |\tilde{c}(\mathtt{Coh})|
\]
\end{proposition}
\begin{proof}
Follows from Lemmas~\ref{lem:catt-term-tag-counts}
and~\ref{lem:frequency-ordering}.  % I may expand this proof in the future.
\end{proof}

% This part and Appendix were created by AI.
\subsection{Optimal Binary Encodings for Composition and Right Whiskering}
We now compute the shortest binary encodings of the programs for $\mathtt{comp}$
and $\mathtt{whiskr}$ in Examples~\ref{ex:comp} and~\ref{ex:whiskr}. We keep the
variable names used in those examples. As above, we use one injective code for
both constructor tags and parameters, and we do not count the length of a
description of the code itself.

The following table gives one code that is optimal for both programs. Its
terminator is $\delta=01$. The last two columns count how often each field
occurs in the respective program.

\begin{center}
\begin{tabular}{lcrr}
\hline
Field & Codeword & $\mathtt{comp}$ & $\mathtt{whiskr}$ \\
\hline
$\mathtt{Var}$       & \texttt{01}     & 11 & 41 \\
$\mathtt{Star}$      & \texttt{001}    &  6 & 20 \\
$\mathtt{ExtendCtx}$ & \texttt{101}    &  5 & 17 \\
$\mathtt{ExtendSub}$ & \texttt{0001}   &  5 & 17 \\
$x$                  & \texttt{1001}   &  5 & 17 \\
$y$                  & \texttt{1101}   &  5 & 17 \\
$z$                  & \texttt{00001}  &  5 & 15 \\
$\mathtt{Arrow}$     & \texttt{10001}  &  3 & 13 \\
$f$                  & \texttt{11001}  &  3 &  9 \\
$g$                  & \texttt{11101}  &  3 &  9 \\
$\mathtt{EmptyCtx}$  & \texttt{000001} &  1 &  3 \\
$\mathtt{EmptySub}$  & \texttt{100001} &  1 &  3 \\
$\mathtt{Coh}$       & \texttt{110001} &  1 &  3 \\
$h$                  & \texttt{111001} &  0 &  5 \\
$a$                  & \texttt{111101} &  0 &  3 \\
\hline
Total field occurrences & & 54 & 192 \\
\hline
\end{tabular}
\end{center}

\paragraph{Example (composition).}
The program for $\mathtt{comp}$ has 33 instructions. Under the displayed code,
the numbers of field occurrences encoded in $2,3,4,5,6$ bits are
$11,11,15,14,3$, respectively. Hence its encoded length is
\[
  2\cdot11+3\cdot11+4\cdot15+5\cdot14+6\cdot3
  =203\text{ bits}.
\]

\paragraph{Example (right whiskering).}
The program for $\mathtt{whiskr}$ has 117 instructions. Under the same code, the
corresponding counts are $41,37,51,46,17$. Hence its encoded length is
\[
  2\cdot41+3\cdot37+4\cdot51+5\cdot46+6\cdot17
  =729\text{ bits}.
\]

\paragraph{Optimality.}
For a fixed terminator, a shortest encoding of a single program assigns the
shortest available words to the most frequent fields: if a more frequent field
has a longer word, exchanging the two words can only shorten the program. Ties
may be settled arbitrarily.

For $\delta=01$, the words of length $\ell\ge2$ in $C_\delta$ are
precisely
\[
  1^i0^{\ell-i-1}1,\qquad 0\le i\le\ell-2.
\]
Thus there is one word of length $2$, two of length $3$, three of length $4$,
and so on. In each example the displayed code gives the shortest available
lengths to the most frequent fields. It is therefore optimal for this
terminator.

It remains to compare terminators. Complementing every bit preserves codeword
lengths, so it suffices to check $0$, $00$, and $01$ among terminators of length
at most two. For $0$, there is one word of each positive length. For $00$, the
fifteen shortest words have lengths
\[
  2,3,4,4,5,5,5,6,6,6,6,6,7,7,7.
\]
Pairing these lengths with the field counts in decreasing order gives
\[
\begin{array}{c|rr}
  \text{Terminator} & \mathtt{comp} & \mathtt{whiskr}\\
  \hline
  0  \text{ or } 1  & 265 & 984\\
  00 \text{ or } 11 & 228 & 820\\
  01 \text{ or } 10 & 203 & 729
\end{array}
\]

Finally, if the terminator has length at least three, there is at most one
admissible word of length three and at most three words of length at most four.
Thus even the three most frequent fields require at least $3,4,4$ bits,
respectively, and every remaining field requires at least $5$ bits in an optimal
frequency ordering. This gives the lower bounds
\[
\begin{aligned}
  \mathtt{comp}:   &\quad 5\cdot54-2\cdot11-6-5=237,\\
  \mathtt{whiskr}: &\quad 5\cdot192-2\cdot41-20-17=841.
\end{aligned}
\]
Both exceed the lengths attained with $01$. The singleton minima are therefore
$203$ and $729$ bits, respectively. Moreover, the optimal worst-case length for
the family $\{\mathtt{comp},\mathtt{whiskr}\}$ is also $729$ bits:
$\mathtt{whiskr}$ alone requires that many bits, and the displayed code attains
the bound for the pair.

\section{Future Work}
\paragraph{Structural Properties and Optimal Encodings}
One may look for structural properties similar to~\ref{lem:catt-term-tag-counts}
and efficient ways of finding worst-case optimal encodings.

\paragraph{Time-Bounded Kolmogorov Complexity}
One can introduce a two-turing machine model for describing time-bounded
Kolmgorov complexity of \catt entities. This model can combine two Turing
machines. The first, $M_1$, takes as input a straight-line program that
constructs a CaTT syntax tree by adding one vertex at a time. Second one, $M_2$,
takes a program containing metaoperations and outputs the expanded straight-line
program used by previous machine. This separation to two turing machines is
particularly useful if we want to limit computation of $M_2$ to expanding
metaoperations to straight-line program of elementary \catt operations, and
$M_1$ to merely outputing encoded syntax tree corresponding to
straight-line program of elementary \catt operations.

\paragraph{Asymptotic Growth}
Given the toolkit for measuring size of \catt entities one may seek asymptotic
bounds of ceratain families of objects.

\bibliography{biblio}




% \section{Asymptotic growth of explicit \catt terms with dimension}
% \section{Growth of explicit \catt terms with dimension}
\appendix
\section{Growth of explicit \catt terms with dimension}
\label{sec:catt-depth-one-minimum}

Disclaimer: This part as well as examples were generated by AI, but I decided to
keep them, at least for now.



Fix \(\Gamma=(x:\obj)\). We ask for the minimum number \(|P|\) of
elementary operations needed to construct an \(n\)-cell in \(\Gamma\).
Here \(n\) is the dimension of its type, with \(\dim(\obj)=0\) and
\(\dim(\arr[A]{s}{t})=\dim(A)+1\). We restrict to terms of coherence
depth at most one~\cite[Section~IV.C]{finster_typetheoretical_2017}:
at each occurrence of \(\coh_{\Delta,A}[\sigma]\), the formal type
\(A\) contains no coherence. Its substitution entries may themselves
be coherences. As before, repeated subtrees are counted separately.

\begin{theorem}[Minimum length at coherence depth one]
\label{thm:catt-depth-one-minimum}
Among such \(n\)-cells, let \(h_n\) be the minimum number of
\(\mathtt{Coh}\) instructions and \(F_n\) the minimum program
length \(|P|\). The \(n\)-fold identity on \(x\) minimizes the
number of occurrences of every tag and hence attains both minima.
For \(n\ge0\),
\[
\begin{aligned}
 h_0&=0, & h_1&=1,  & h_{n+2}&=2h_{n+1}+h_n,\\
 F_0&=1, & F_1&=11, & F_{n+2}&=2F_{n+1}+F_n+6n+12.
\end{aligned}
\]
In particular, \(h_n\) is the \(n\)th Pell
number~\cite[Chapter~7]{koshy_pell_2014}.
\end{theorem}

\begin{proof}
Put \(t_0=x\) and \(t_{k+1}=\mathrm{id}_{t_k}\), using the disc
identity of~\cite[Definitions~2.3--2.4]{Beyond_eh}. Its formal
type contains no coherence. Written in full, its substitution
contains \(t_k\) once and \(t_j\) twice for each \(j<k\), once
for each of the two \(j\)-dimensional faces. Thus all \(t_k\)
have coherence depth at most one.

We prove by strong induction that \(t_k\) minimizes every tag
count. The only \(0\)-cell in \(\Gamma\) is \(x\). Any competing
\((k+1)\)-cell is of the form \(\coh_{\Delta,A}[\sigma]\).
Since \(A\) is coherence-free, its source is a \(k\)-dimensional
variable. This variable and its faces determine a map from the
\(k\)-disc into the pasting scheme of \(\Delta\), injective
by~\cite[Lemma~29]{finster_typetheoretical_2017}. Each disc
variable therefore selects a distinct entry of \(\sigma\).
By induction, each selected entry contains at least as many
occurrences of every tag as the corresponding \(t_j\), and the
explicit tree counts these entries separately. The selected
declaration types and the formal type have the same tag counts
as their counterparts in \(t_{k+1}\): a coherence-free \(j\)-type
has one \(\mathtt{Star}\), \(j\) arrows and \(2j\) variable
endpoints. Both lists have one extension per declaration and one
empty node. Adding the coherence root proves the claim.

It remains to count the \(t_k\). The initial pairs \((h_k,F_k)\)
are \((0,1)\) and \((1,11)\): for \(t_1\), the root, context,
formal type and substitution contribute \(1+3+4+3\) vertices.
Compared with \(D_n\), the disc \(D_{n+1}\) has one additional
declaration in each of dimensions \(n\) and \(n+1\). Hence the
substitution for \(t_{n+2}\) gains one copy each of \(t_n\) and
\(t_{n+1}\) over that for \(t_{n+1}\). This gives the recurrence
for \(h_n\). The two declaration types add \(3n+1\) and \(3n+4\)
vertices, the two lists add four extensions, and the formal type
adds one arrow and two endpoints. Their sum is \(6n+12\), giving
the recurrence for \(F_n\).
\end{proof}

\begin{corollary}[Growth and cost of the tag fields]
The minimum program length is
\[
 F_n=7h_{n+1}+6h_n-3n-6.
\]
Both \(h_n\) and \(F_n\) grow as \(\Theta((1+\sqrt{2})^n)\).
For any fixed code \(c\), the \(n\)-fold identity also minimizes
the total length of the tag fields.
\end{corollary}

\begin{proof}
Both \(F_n+3n+6\) and \(7h_{n+1}+6h_n\) satisfy the Pell recurrence
with initial values \(7,20\), so they agree. The growth rate follows
from the Pell formula
\(h_n=((1+\sqrt{2})^n-(1-\sqrt{2})^n)/(2\sqrt{2})\).
For the last claim, multiply each tag count by its codeword length
and sum over the eight tags. This counts only the tag fields;
\(\|P\|\) also includes the parameter fields.
\end{proof}

For example, the first six minimum program lengths are
\(1,11,35,99,257,643\). In dimension two the minimum is \(35\)
operations, of which two are \(\mathtt{Coh}\) instructions.

\end{document}

%%% Local Variables:
%%% mode: LaTeX
%%% TeX-master: "main"
%%% End:
